% !TeX root = ../eccv2016submission.tex

Convolutional Neural Networks (CNNs) were arguably popularized within the vision community in 2009 through AlexNet~\cite{krizhevsky2009learning} and its celebrated victory at the ImageNet competition~\cite{deng2009imagenet}.
Since then there has been a notable shift towards CNNs in many areas of computer vision~\cite{krizhevsky2012imagenet,sermanet2013overfeat,simonyan2014very,springenberg2014striving,szegedy2015going,he2015deep}. As this shift unfolds, a second trend emerges; deeper and deeper CNN architectures are being developed and trained. Whereas AlexNet had 5 convolutional layers~\cite{krizhevsky2009learning}, the VGG network and GoogLeNet in 2014 had 19 and 22 layers respectively~\cite{simonyan2014very,szegedy2015going}, and most recently the ResNet architecture featured 152 layers~\cite{he2015deep}.

Network depth is a major determinant of model expressiveness, both in theory \cite{haastad1991power,haastad1987computational} and in practice \cite{he2015deep,simonyan2014very,szegedy2015going}. However, very deep models also introduce new challenges: vanishing gradients in backward propagation, diminishing feature reuse in forward propagation, and long training time.

\emph{Vanishing Gradients} is a well known nuisance in neural networks with many layers~\cite{bengio1994learning}. As the gradient information is back-propagated, repeated multiplication or convolution with small weights renders the gradient information ineffectively small in earlier layers. Several approaches exist to reduce this effect in practice, for example through careful initialization \cite{glorot2010understanding}, hidden layer supervision \cite{lee2014deeply}, or, recently, Batch Normalization~\cite{ioffe2015batch}.

\emph{Diminishing feature reuse} during forward propagation (also known as loss in information flow \cite{srivastava2015highway}) refers to the analogous problem to vanishing gradients in the forward direction. The features of the input instance, or those computed by earlier layers, are ``washed out'' through repeated multiplication or convolution with (randomly initialized) weight matrices, making it hard for later layers to identify and learn ``meaningful'' gradient directions. Recently, several new architectures attempt to circumvent this problem through direct identity mappings between layers, which allow the network to pass on features unimpededly from earlier layers to later layers~\cite{he2015deep,srivastava2015highway}.

\emph{Long training time} is a serious concern as networks become very deep. The forward and backward passes scale linearly with the depth of the network. Even on modern computers with multiple state-of-the-art GPUs, architectures like the 152-layer ResNet require several weeks to converge on the ImageNet dataset~\cite{he2015deep}.

The researcher is faced with an inherent dilemma:
shorter networks have the advantage that information flows efficiently forward and backward, and can therefore be trained effectively and within a reasonable amount of time. However, they are not expressive enough to represent the complex concepts that are commonplace in computer vision applications. Very deep networks have much greather model complexity, but are very difficult to train in practice and require a lot of time and patience.

In this paper, we propose \emph{deep networks with \name{}}, a novel training algorithm that is based on the seemingly contradictory insight that ideally we would like to have a \emph{deep} network during \emph{testing} but a \emph{short} network during \emph{training}. We resolve this conflict by creating deep Residual Network~\cite{he2015deep} architectures (with hundreds or even thousands of layers) with sufficient modeling capacity; however, during training  we shorten the network significantly by randomly removing a substantial fraction of layers independently for each sample or mini-batch.
The effect is a network with a small \emph{expected} depth during training, but a large depth during testing. Although seemingly simple, this approach is surprisingly effective in practice.

In extensive experiments we observe that training with \name{} substantially reduces training time and test error (resulting in multiple new records to the best of our knowledge at the time of initial submission to ECCV).
The reduction in training time can be attributed to the shorter forward and backward propagation, so the training time no longer scales with the full depth, but the shorter \emph{expected depth} of the network. We attribute the reduction in test error to two factors: 1) shortening the (expected) depth during training reduces the chain of forward propagation steps and gradient computations, which strengthens the gradients especially in earlier layers during backward propagation; 2) networks trained with \name{} can be interpreted as an implicit \emph{ensemble} of networks of different depths, mimicking the record breaking ensemble of depth varying ResNets trained by He et al.~\cite{he2015deep}.

We also observe that similar to Dropout \cite{srivastava2014dropout}, training with \name{} acts as a regularizer, even in the presence of Batch Normalization \cite{ioffe2015batch}. On experiments with CIFAR-10, we increase the depth of a ResNet beyond 1000 layers and still obtain significant improvements in test error.
